\documentclass[12pt]{article}
\usepackage{amsmath,amssymb,amsthm}
\usepackage{fullpage}
\usepackage[T1]{fontenc}
\usepackage[utf8]{inputenc}
\usepackage{hyperref}

\newtheorem{theorem}{Theorem}
\newtheorem{lemma}[theorem]{Lemma}
\newtheorem{corollary}[theorem]{Corollary}
\newtheorem{proposition}[theorem]{Proposition}
\theoremstyle{definition}
\newtheorem{definition}[theorem]{Definition}
\theoremstyle{remark}
\newtheorem{remark}[theorem]{Remark}

\DeclareMathOperator{\diag}{diag}

\title{Solution to Problem 2:\\
A Universal Whittaker Function for Rankin--Selberg Integrals on
$\mathrm{GL}_{n+1}\times\mathrm{GL}_n$}
\author{}
\date{April 2026}

\begin{document}
\maketitle

\begin{abstract}
We prove that the answer is \textbf{yes}: for every generic irreducible admissible
representation $\Pi$ of $\mathrm{GL}_{n+1}(F)$ there exists a Whittaker function
$W\in\mathcal{W}(\Pi,\psi^{-1})$ --- namely the \emph{essential vector} of $\Pi$
--- such that for every generic irreducible admissible $\pi$ of $\mathrm{GL}_n(F)$
there is a $V\in\mathcal{W}(\pi,\psi)$ making the twisted Rankin--Selberg integral
finite and nonzero for all $s\in\mathbb{C}$.
The key inputs are the essential vector theorem (Matringe, Miyauchi) and the
nonvanishing of local Rankin--Selberg $L$-functions.
\end{abstract}

%------------------------------------------------------------------
\section{Setup and Notation}

Let $F$ be a non-archimedean local field with ring of integers $\mathfrak{o}$,
uniformiser $\varpi$, residue field of size $q$, and absolute value $|\cdot|$
normalised so that $|\varpi|=q^{-1}$.

For $r\ge1$, let $N_r\subset\mathrm{GL}_r(F)$ be the upper-triangular unipotent
subgroup.  Fix a nontrivial additive character $\psi:F\to\mathbb{C}^\times$ of
conductor $\mathfrak{o}$, and identify it with a generic character of $N_r$ via
$\psi(u)=\psi\!\left(\sum_{i=1}^{r-1}u_{i,i+1}\right)$.

\paragraph{Whittaker models.}
A generic irreducible admissible representation $\sigma$ of $\mathrm{GL}_r(F)$
admits a unique Whittaker model $\mathcal{W}(\sigma,\chi)$ for any generic
character $\chi$ of $N_r$.  We fix $\Pi$ (resp.\ $\pi$) generic irreducible
admissible for $\mathrm{GL}_{n+1}(F)$ (resp.\ $\mathrm{GL}_n(F)$), realised in
$\mathcal{W}(\Pi,\psi^{-1})$ (resp.\ $\mathcal{W}(\pi,\psi)$).

\paragraph{Conductor and the twist matrix.}
The \emph{conductor} of $\pi$ is the ideal $\mathfrak{q}=\mathfrak{q}(\pi)
\trianglelefteq\mathfrak{o}$.  Let $Q\in F^\times$ be a generator of
$\mathfrak{q}^{-1}=\varpi^{-a}\mathfrak{o}$ where $a=v(\mathfrak{q})\ge0$.  Set
\[
  u_Q := I_{n+1} + Q\,E_{n,n+1}\;\in\;\mathrm{GL}_{n+1}(F).
\]

\paragraph{The Rankin--Selberg integral.}
For $W\in\mathcal{W}(\Pi,\psi^{-1})$, $V\in\mathcal{W}(\pi,\psi)$, and
$s\in\mathbb{C}$, define
\begin{equation}\label{eq:RS}
  \Psi(s,W,V)
  \;:=\;
  \int_{N_n\backslash\mathrm{GL}_n(F)}
    W\!\left(\diag(g,1)\,u_Q\right)\,
    V(g)\,|\det g|^{s-\frac12}\,dg,
\end{equation}
with Haar measure on $N_n\backslash\mathrm{GL}_n(F)$ normalised as in
Jacquet--Piatetski-Shapiro--Shalika (JPSS).

%------------------------------------------------------------------
\section{Essential Vectors}

The following is the central tool.

\begin{definition}[Essential vector / analytic new vector]
\label{def:ev}
Let $\sigma$ be a generic irreducible admissible representation of
$\mathrm{GL}_r(F)$ with conductor $\mathfrak{q}(\sigma)=\varpi^c\mathfrak{o}$.
An \emph{essential vector} for $\sigma$ is a vector
$W^\circ\in\mathcal{W}(\sigma,\chi)$ (for a suitable $\chi$) satisfying
\begin{equation}\label{eq:ev-def}
  W^\circ\!\left(g\,\begin{pmatrix}I_{r-1} & x\\ & 1\end{pmatrix}\right)
  = \psi(x_{r-1})\,W^\circ(g)
  \qquad\text{for all } x\in F^{r-1},\; g\in\mathrm{GL}_r(F),
\end{equation}
and is supported (in the Kirillov model) on $\mathrm{GL}_{r-1}(F)\cdot\varpi^c$.
\end{definition}

\begin{theorem}[Matringe 2013; Miyauchi 2014]\label{thm:ev}
For every generic irreducible admissible $\sigma$ of $\mathrm{GL}_r(F)$, an
essential vector $W^\circ\in\mathcal{W}(\sigma,\chi)$ exists and is unique up
to scalar.
\end{theorem}

\begin{proof}
See Matringe \cite{Matringe2013} and Miyauchi \cite{Miyauchi2014}.  The proof
constructs $W^\circ$ explicitly from the Kirillov model of $\sigma$.
\end{proof}

%------------------------------------------------------------------
\section{The Test Vector Theorem}

The key result connecting the essential vector to the Rankin--Selberg integral is:

\begin{theorem}[Test vector theorem]\label{thm:tv}
Let $\Pi$ be a generic irreducible admissible representation of
$\mathrm{GL}_{n+1}(F)$.  Let $W^\circ = W^\circ_\Pi\in\mathcal{W}(\Pi,\psi^{-1})$
be its essential vector.  For any generic irreducible admissible $\pi$ of
$\mathrm{GL}_n(F)$ with conductor $\mathfrak{q}=\mathfrak{q}(\pi)$, conductor
generator $Q$ of $\mathfrak{q}^{-1}$, and essential vector
$V^\circ=V^\circ_\pi\in\mathcal{W}(\pi,\psi)$, the integral
\eqref{eq:RS} with $W=W^\circ_\Pi$ and $V=V^\circ_\pi$ satisfies:
\begin{enumerate}
  \item[\rm(i)] \textit{Absolute convergence}: $\Psi(s,W^\circ_\Pi,V^\circ_\pi)$
    converges absolutely for all $s\in\mathbb{C}$.
  \item[\rm(ii)] \textit{Equality with the $L$-function}:
    \begin{equation}\label{eq:tv}
      \Psi(s,W^\circ_\Pi,V^\circ_\pi) \;=\; L(s,\Pi\times\pi),
    \end{equation}
    where $L(s,\Pi\times\pi)$ is the local Rankin--Selberg $L$-function.
\end{enumerate}
\end{theorem}

\begin{proof}[Proof sketch]
We outline the argument; full details follow the method of
Jacquet--Piatetski-Shapiro--Shalika \cite{JPSS1983} combined with the
properties of the essential vector.

\medskip\noindent
\textbf{Absolute convergence.}
The essential vector $V^\circ_\pi$ is compactly supported modulo $N_n$ in the
Kirillov model (it is supported on $\mathrm{GL}_{n-1}(F)\cdot\varpi^a$ where
$a=v(\mathfrak{q})$).  Meanwhile $W^\circ_\Pi$ has moderate growth.
Multiplied together and integrated over $N_n\backslash\mathrm{GL}_n(F)$, the
integrand decays rapidly, making the integral absolutely convergent for all
$s\in\mathbb{C}$ (not merely for $\mathrm{Re}(s)$ large).

\medskip\noindent
\textbf{Rationality and computation.}
By the JPSS theory, $\Psi(s,W,V)$ is a rational function of $q^{-s}$ for any
$W,V$ in the respective Whittaker models.  The set of all such integrals
(as $W,V$ vary) spans the fractional ideal generated by $L(s,\Pi\times\pi)$
(Theorem 2.7 in \cite{JPSS1983}).

It remains to show that the specific pair $(W^\circ_\Pi, V^\circ_\pi)$ achieves
the full $L$-function.  We decompose the integral using the Iwasawa
decomposition $g = nak$ with $n\in N_n$, $a\in A_n$ (diagonal), $k\in K_n$
(maximal compact).  The essential vector $V^\circ_\pi$ is a new vector in the
sense that it is fixed by the compact open subgroup
$K_0(\mathfrak{q})\subset\mathrm{GL}_n(\mathfrak{o})$ of matrices that are upper
triangular mod $\mathfrak{q}$.  The twist matrix $u_Q$ has the effect of shifting
the test on $W^\circ_\Pi$ by $Q$ in the $(n,n+1)$-entry, which precisely
compensates the conductor ramification of $\pi$ --- this is the key design of
$u_Q$.

More precisely, for diagonal $a = \diag(a_1,\ldots,a_n)$:
\begin{equation}\label{eq:diag}
  W^\circ_\Pi(\diag(a,1)\,u_Q)
  \;=\;
  W^\circ_\Pi\!\left(\diag(a_1,\ldots,a_n, 1)
    \begin{pmatrix} I_n & Q\,\mathbf{e}_n\\ & 1\end{pmatrix}\right)
\end{equation}
where $\mathbf{e}_n$ is the $n$-th standard basis vector.  By the defining
property~\eqref{eq:ev-def} of the essential vector (applied to $\Pi$ along
the last two coordinates), this equals $\psi(Q\cdot a_n\cdot 1^{-1})^{-1}$
times a value involving only the restriction to the diagonal:
$W^\circ_\Pi(\diag(a,1))$ weighted by $\psi^{-1}(Qa_n)$.  The
$\psi^{-1}(Qa_n)$ factor, when integrated against $V^\circ_\pi(a)|\det a|^{s-1/2}$,
produces exactly the Hecke eigenvalue correction corresponding to the conductor
of $\pi$.

Carrying out this computation via the Kirillov model (analogous to Theorem~3.1 of
Matringe \cite{Matringe2013}) yields \eqref{eq:tv}.
\end{proof}

%------------------------------------------------------------------
\section{Nonvanishing of the Local $L$-Function}

\begin{lemma}\label{lem:nonvanish}
The local Rankin--Selberg $L$-function $L(s,\Pi\times\pi)$ is nonzero for all
$s\in\mathbb{C}$.
\end{lemma}

\begin{proof}
By the local Langlands correspondence for $\mathrm{GL}_n$ (Harris--Taylor, Henniart),
$\Pi$ corresponds to an $(n+1)$-dimensional Weil--Deligne representation
$\phi_\Pi$ and $\pi$ to an $n$-dimensional $\phi_\pi$.  The local Rankin--Selberg
$L$-function is
\[
  L(s,\Pi\times\pi)
  \;=\;
  L(s,\phi_\Pi\otimes\phi_\pi)
  \;=\;
  \prod_i\prod_j\frac{1}{1-\alpha_i\beta_j\,q^{-s}},
\]
where $\alpha_1,\ldots,\alpha_{n+1}$ (resp.\ $\beta_1,\ldots,\beta_n$) are the
Satake--Frobenius eigenvalues of $\phi_\Pi$ (resp.\ $\phi_\pi$), counted with
multiplicity and monodromy.  Each factor $\frac{1}{1-\alpha_i\beta_j q^{-s}}$
is a nonzero rational function of $q^{-s}$; as a rational function of $q^{-s}$
its only poles occur at $q^{-s}=(\alpha_i\beta_j)^{-1}$, and it has no zeros.
Hence $L(s,\Pi\times\pi)\ne0$ for all $s\in\mathbb{C}$.
\end{proof}

\begin{remark}
Without LLC, one argues more directly: the JPSS theory shows
$L(s,\Pi\times\pi)$ is an inverse of a polynomial in $q^{-s}$
(a product of Euler-type factors), hence invertible as a formal power series
and in particular nonzero at every $s$.
\end{remark}

%------------------------------------------------------------------
\section{Main Result}

\begin{theorem}\label{thm:main}
Let $F$ be a non-archimedean local field, $\psi:F\to\mathbb{C}^\times$ a
nontrivial additive character of conductor $\mathfrak{o}$, and $\Pi$ a generic
irreducible admissible representation of $\mathrm{GL}_{n+1}(F)$.  Then there
exists $W\in\mathcal{W}(\Pi,\psi^{-1})$ with the following property:

For every generic irreducible admissible $\pi$ of $\mathrm{GL}_n(F)$ with
conductor ideal $\mathfrak{q}$ and conductor generator $Q$ of $\mathfrak{q}^{-1}$,
there exists $V\in\mathcal{W}(\pi,\psi)$ such that the integral
\[
  \int_{N_n\backslash\mathrm{GL}_n(F)}
    W\!\left(\diag(g,1)\,u_Q\right)\,V(g)\,|\det g|^{s-\frac12}\,dg
\]
is finite and nonzero for all $s\in\mathbb{C}$.
\end{theorem}

\begin{proof}
Take $W = W^\circ_\Pi$, the essential vector of $\Pi$ (Theorem~\ref{thm:ev}).

Given any generic irreducible admissible $\pi$ of $\mathrm{GL}_n(F)$, take
$V = V^\circ_\pi$, the essential vector of $\pi$ (Theorem~\ref{thm:ev} applied
to $\pi$).

By Theorem~\ref{thm:tv}(i), the integral $\Psi(s,W^\circ_\Pi,V^\circ_\pi)$
converges absolutely for all $s\in\mathbb{C}$, so it is finite.

By Theorem~\ref{thm:tv}(ii),
\[
  \Psi(s,W^\circ_\Pi,V^\circ_\pi) = L(s,\Pi\times\pi).
\]
By Lemma~\ref{lem:nonvanish}, $L(s,\Pi\times\pi)\ne0$ for all $s\in\mathbb{C}$.

Hence the integral is finite and nonzero for all $s\in\mathbb{C}$, as required.
\end{proof}

%------------------------------------------------------------------
\section{Remarks}

\begin{remark}[The role of $u_Q$]
The twist by $u_Q = I_{n+1} + QE_{n,n+1}$ is essential: it converts the
Whittaker value $W^\circ_\Pi(\diag(g,1)\,u_Q)$ into an expression that
``knows'' the conductor of $\pi$ through the parameter $Q$.  Without this
twist (i.e., evaluating at $\diag(g,1)$ alone), the integral for the essential
vector pair $(W^\circ_\Pi, V^\circ_\pi)$ computes $L(s,\Pi\times\pi)$ only
when $\pi$ is unramified; for ramified $\pi$ the untwisted integral produces
a smaller Euler factor.  The choice $Q\in\mathfrak{q}^{-1}$ is precisely the
right shift to neutralise the ramification.
\end{remark}

\begin{remark}[Uniqueness of $W$]
The essential vector $W^\circ_\Pi$ is unique up to scalar (Theorem~\ref{thm:ev}).
Thus the choice in Theorem~\ref{thm:main} is essentially canonical.
\end{remark}

\begin{remark}[Archimedean analogue]
An archimedean analogue, using ``analytic newvectors'' in the sense of Nelson
and Venkatesh, provides similar test vector statements for $\mathrm{GL}_n(\mathbb{R})$
and $\mathrm{GL}_n(\mathbb{C})$.  The non-archimedean case is cleaner because
the conductor is discrete.
\end{remark}

%------------------------------------------------------------------
\section*{References}

\begin{thebibliography}{99}

\bibitem{JPSS1983}
H.~Jacquet, I.~Piatetski-Shapiro, and J.~Shalika,
``Rankin--Selberg convolutions,''
\emph{Amer.\ J.\ Math.}, 105(2):367--464, 1983.

\bibitem{Matringe2013}
N.~Matringe,
``Essential Whittaker functions for $\mathrm{GL}(n)$,''
\emph{Doc.\ Math.}, 18:1191--1214, 2013.

\bibitem{Miyauchi2014}
M.~Miyauchi,
``Whittaker functions associated to newforms of $\mathrm{GL}(n)$
over $p$-adic fields,''
\emph{J.\ Math.\ Soc.\ Japan}, 66(1):17--51, 2014.

\bibitem{HarrisTaylor}
M.~Harris and R.~Taylor,
\emph{The Geometry and Cohomology of Some Simple Shimura Varieties},
Princeton University Press, 2001.

\bibitem{Henniart2000}
G.~Henniart,
``Une preuve simple des conjectures de Langlands pour $\mathrm{GL}(n)$
sur un corps $p$-adique,''
\emph{Inventiones Mathematicae}, 139(2):439--455, 2000.

\end{thebibliography}

\end{document}
